\documentclass[10pt,a4paper]{amsart}
\usepackage[margin=19mm]{geometry}
\usepackage{amsmath,amssymb,mathtools}
\usepackage[scr=rsfs,cal=euler]{mathalfa}
\usepackage{xfrac}
\usepackage[unicode,hidelinks,bookmarks=false]{hyperref}
\newcommand{\ie}{i.e.\ }
\newcommand{\cf}{cf.\ }
\newcommand{\resp}{respectively\ }
\newcommand{\Subsection}{\subsection}
\providecommand{\nobreakdash}{\nobreak\hbox{-}}
\setlength{\emergencystretch}{2em}
\multlinegap0pt
\usepackage{bm}
\usepackage{bbm}
\usepackage{dsfont}
\usepackage{xspace}
\usepackage{cancel}
\usepackage{mathtools}
\usepackage{amsthm}
\makeatletter
\@ifundefined{Thm}{\newtheorem{Thm}{Thm}}{}
\@ifundefined{Prop}{\newtheorem{Prop}[Thm]{Proposition}}{}
\makeatother
\newcommand{\br}{\medskip\noindent}
\newcommand{\ip}[2]{\left \langle{#1},{#2} \right \rangle}
\newcommand{\mc}[1]{\mathcal{#1}}
\newcommand{\set}[1]{\left\{#1\right\}}
\title[Some Fock spaces with depth two action]{Some Fock spaces with depth two action}
\author{Michael Anshelevich, Jacob Mashburn}
\date{Source: arXiv:2103.13936v2 (CC BY 4.0)}
\begin{document}
\maketitle

\noindent\emph{Source and licence.} Some Fock spaces with depth two action. Version arXiv:2103.13936v2 (CC BY 4.0), distributed under the Creative Commons Attribution 4.0 International licence, as recorded in the arXiv licence metadata for this exact version and shown on the abstract page https://arxiv.org/abs/2103.13936v2. Full rights notice: licenses/arxiv-2103.13936-CC-BY-4.0.txt.\par\smallskip

\noindent\textit{Editorial scope.} Counted unit: the Prop whose source label is \texttt{Prop:Interacting-traciality} in the article (source0.tex lines 1025--1043, source label \texttt{Prop:Interacting-traciality}), delivered together with the article's own complete original proof of it (source0.tex lines 1046--1128, one \texttt{proof} environment of 83 lines). No mathematics is written, repaired or re-derived: statement and proof are the article's own text line for line. Required context retained uncounted: none. Every definition, display and label of those lines is retained unchanged. Omitted from this package and not redistributed: 1--1024, 1044--1045, 1129--1533. Nothing inside a retained excerpt is omitted or altered: every retained body is delivered exactly as the article's lines are written, so no omission receipt of any kind accompanies this package. Citations: the retained excerpts carry no citation command at all, so no bibliography entry of the article is transcribed and no reference is redistributed. Reference adaptation: none was needed and none was made; every reference inside the delivered unit resolves inside it, so no unresolved reference or citation is produced. Printed slips kept literal: no printed word, symbol, hypothesis or number of the counted statement or of its proof was altered. Layout and class: the article's own class is \texttt{amsart} with packages including times, amssymb, comment, color, hyperref; the wrapper is the lane's pinned \texttt{amsart} editorial wrapper at 10pt on A4, which restates the theorem-like environments the retained text needs and adds the article's own macro definitions that the retained text needs (\texttt{\textbackslash br}, \texttt{\textbackslash ip}, \texttt{\textbackslash mc}, \texttt{\textbackslash set}), with extra wrapper packages (bm, bbm, dsfont, xspace, cancel, mathtools). The delivered numbering is the unit's own. Assets: no retained span contains a figure environment, a graphics-inclusion command, a PDF-page inclusion, a table or a TikZ picture; no image file is copied into this package, the package declares no asset dependency and no image file is redistributed with it. Licence: version arXiv:2103.13936v2 of this article is distributed under the Creative Commons Attribution 4.0 International licence, as recorded in the arXiv licence metadata for this exact version and shown on the abstract page https://arxiv.org/abs/2103.13936v2. A full rights notice is delivered at licenses/arxiv-2103.13936-CC-BY-4.0.txt. Characters: every retained excerpt is pure ASCII, so the delivered engine, XeLaTeX with its default Latin Modern text font, reports no missing character.
\begin{Prop}
\label{Prop:Interacting-traciality}
\

\begin{enumerate}
\item
Suppose that $\Gamma^{alg}_{\gamma, \Lambda}(\mc{B}, \phi)$ and $\Gamma^{alg}_{\gamma, \Lambda}(\mc{B}, \phi; r)$ commute. Then $W$ extends to a linear map on $\mc{F}_{\text{alg}}(\mc{B})$,
\[
\Gamma^{alg}_{\gamma, \Lambda}(\mc{B}, \phi) = \set{W(\vec{\xi}) : \vec{\xi} \in \mc{F}_{\text{alg}}(\mc{B})}, \quad \Gamma^{alg}_{\gamma, \Lambda}(\mc{B}, \phi; r) = \set{W_r(\vec{\xi}) : \vec{\xi} \in \mc{F}_{\text{alg}}(\mc{B})},
\]
and for $\vec{\xi} \in \mc{F}_{\text{alg}}(\mc{B})$,
\[
W(\vec{\xi})^\ast = W(S(\vec{\xi})),
\]
Therefore for $A \in \Gamma^{alg}_{\gamma, \Lambda}(\mc{B}, \phi)$, $S(A \Omega) = A^\ast \Omega$.
\item
In addition to the assumption in (a), suppose that $\phi$ is tracial on $\mc{B}$. Then the vacuum state is tracial on $\Gamma^{alg}_{\gamma, \Lambda}(\mc{B}, \phi)$.
\end{enumerate}
\end{Prop}

\begin{proof}
\

\begin{enumerate}
\item
Since $\Omega$ is cyclic for $\Gamma^{alg}_{\gamma, \Lambda}(\mc{B}, \phi)$ and $\Gamma^{alg}_{\gamma, \Lambda}(\mc{B}, \phi; r)$, it is separating for them (and the von Neumann algebras they generate). Using the recursion and induction, for $\vec{\xi} \in \mc{F}_{\text{alg}}(\mc{B})$
\[
\begin{split}
W(a^+(b)(\vec{\xi})) \Omega
& = X(b) W(\vec{\xi}) \Omega - W(a^0(b) (\vec{\xi})) \Omega - W(a^-(b)(\vec{\xi})) \Omega \\
& = X(b) \vec{\xi} - a^0(b) (\vec{\xi}) - a^-(b)(\vec{\xi}) \\
& = a^+(b)(\vec{\xi}).
\end{split}
\]
In particular, if $\vec{\xi} = 0$, then $W(\vec{\xi}) \Omega = 0$, so  $W(\vec{\xi}) = 0$, and the map $W$ is well defined.
\\
Clearly, $W(\vec{\xi}) \in \Gamma^{alg}_{\gamma, \Lambda}(\mc{B}, \phi)$. On the other hand, using induction on the number of factors, a monomial $X(u_0) X(u_1) \ldots X(u_n)$ is a linear combination of terms of the form $X(u_0) W(\vec{\xi})$, each of which is a linear combination of the terms of the form $W(\vec{\xi'})$ by the recursion.


Next, note that
\[
X_r(b) \Omega = X(b) \Omega
\]
and
\[
W(b)^\ast = X(b)^\ast = X(b^\ast) = W(b^\ast).
\]
Using induction,
\[
\begin{split}
W(a^+(b)(\vec{\xi}))^\ast
& = W(\vec{\xi})^\ast X(b)^\ast - W(a^0(b) (\vec{\xi}))^\ast - W(a^-(b)(\vec{\xi}))^\ast \\
& = W(S(\vec{\xi})) X(b^\ast) - W(S(a^0(b) (\vec{\xi}))) - W(S(a^-(b)(\vec{\xi}))).
\end{split}
\]
On the other hand, if $X_r(b)$ commutes with $\Gamma^{alg}_{\gamma, \Lambda}(\mc{B}, \phi)$,
\[
\begin{split}
W(S(\vec{\xi})) X(b^\ast) \Omega
& = W(S(\vec{\xi})) X_r(b^\ast) \Omega
= X_r(b^\ast) W(S(\vec{\xi})) \Omega \\
& = S X(b) S S(\vec{\xi})
= S X(b) (\vec{\xi}).
\end{split}
\]
Since $\Omega$ is separating for $\Gamma^{alg}_{\gamma, \Lambda}(\mc{B}, \phi)$, this implies that
\[
W(S(\vec{\xi})) X(b^\ast)
= W(S X(b) (\vec{\xi}))
= W(S(a^+(b)(\vec{\xi}))) + W(S(a^0(b)(\vec{\xi}))) + W(S(a^-(b)(\vec{\xi}))).
\]
Thus
\[
W(a^+(b)(\vec{\xi}))^\ast = W(S(a^+(b)(\vec{\xi}))).
\]
\item
We first show that $X_r(b)^\ast \Omega = X(b)^\ast \Omega$. Indeed, for $\vec{\xi} \in \mc{B}^{\otimes n}$, $\ip{X_r(b) (\vec{\xi})}{\Omega} = 0$ if $n \neq 1$. For $u \in \mc{B}$, using the fact that $\phi$ is a trace,
\[
\begin{split}
\ip{X_r(b) (u)}{\Omega}
& = \ip{S a^-(b^\ast)(S(u))}{\Omega}
= \overline{\phi[b^\ast u]}
= \phi[u b]
= \phi[b u]
= \ip{u}{b^\ast}
= \ip{u}{X(b)^\ast \Omega}
\end{split}
\]

\br
To check that the vacuum state is tracial on $\Gamma^{alg}_{\gamma, \Lambda}(\mc{B}, \phi)$, it suffices to verify that
\[
\begin{split}
\ip{W(\vec{\xi}) X(b) \Omega}{\Omega}
& = \ip{W(\vec{\xi}) X_r(b) \Omega}{\Omega}
= \ip{X_r(b) W(\vec{\xi}) \Omega}{\Omega} \\
& = \ip{W(\vec{\xi}) \Omega}{X_r(b)^\ast \Omega}
= \ip{W(\vec{\xi}) \Omega}{X(b)^\ast \Omega} \\
& = \ip{X(b) W(\vec{\xi}) \Omega}{\Omega}. \qedhere
\end{split}
\]
\end{enumerate}
\end{proof}

\end{document}
